\section*{Chapter \ref{ch:g12:classes-and-objects} --- Classes and Objects}

\begin{solution}{exo:g12:classes-and-objects:1}
Something like \code{<__main__.Point object at 0x7f3a...>}: the \omterm{def:g12:classes-and-objects:class}{class} name and a memory address, which changes from
run to run. Defining a \omterm{def:g12:classes-and-objects:special-method}{special method} gives a readable form:

\smallskip\noindent\hspace*{2em}\code{def __repr__(self):}

\noindent\hspace*{4em}\code{return f"({self.x}, {self.y})"}

\smallskip\noindent Without \code{__str__}, \code{print} uses \code{__repr__}.
\end{solution}

\begin{solution}{exo:g12:classes-and-objects:2}
$b = \frac12 + \frac14 = \frac34$ and $c = \frac34 \times \frac23 = \frac{6}{12} = \frac12$: \code{print(b, c)} shows
\code{3/4 1/2}, and \code{print(repr(c))} shows \code{Fraction(1, 2)}.
\end{solution}

\begin{solution}{exo:g12:classes-and-objects:3}
\smallskip\noindent\hspace*{2em}\code{class Tally:}

\noindent\hspace*{4em}\code{def __init__(self):}

\noindent\hspace*{6em}\code{self._count = 0}

\noindent\hspace*{4em}\code{def add(self, k=1):}

\noindent\hspace*{6em}\code{self._count += k}

\noindent\hspace*{4em}\code{def value(self):}

\noindent\hspace*{6em}\code{return self._count}
\end{solution}

\begin{solution}{exo:g12:classes-and-objects:4}
\code{p} and \code{q} are one object, moved once: both are at $(3, 2)$. \code{r} is a new object built from those
values, then moved: $(3, 3)$; moving it changed neither \code{p} nor \code{q}.
\end{solution}

\begin{solution}{exo:g12:classes-and-objects:5}
\code{p.move(1, 1)} is the call \code{Point.move(p, 1, 1)}: three arguments, the object and the two numbers. The
\omterm{def:g12:classes-and-objects:class}{method} was written with two parameters, as if Python did not pass the object. The repair is
\code{def move(self, dx, dy):}; without it, even the body could not name the object it should change.
\end{solution}

\begin{solution}{exo:g12:classes-and-objects:6}
\smallskip\noindent\hspace*{2em}\code{def sub(a, b):}

\noindent\hspace*{4em}\code{num = a.numerator() * b.denominator() - b.numerator() * a.denominator()}

\noindent\hspace*{4em}\code{return Fraction(num, a.denominator() * b.denominator())}

\smallskip\noindent The result is built by the \omterm{def:g12:classes-and-objects:constructor}{constructor}, which reduces it: the invariant is the \omterm{def:g12:classes-and-objects:class}{class}'s work, not
the user's. \code{sub(Fraction(1, 2), Fraction(1, 3))} is \code{1/6}.
\end{solution}

\begin{solution}{exo:g12:classes-and-objects:7}
A set finds its elements by their hash, and equal elements must have equal hashes, or two equal fractions could sit
in two places of the same set. The hash Python gives every object by default comes from its identity: two equal
fractions, being two objects, would get different hashes. So a \omterm{def:g12:classes-and-objects:class}{class} that defines \code{__eq__} loses that default,
and its objects are refused by sets and as dictionary keys until it defines \code{__hash__} from the values it
compares, for example \code{return hash((self._num, self._den))}. This is only safe because no \omterm{def:g12:classes-and-objects:class}{method} changes a
fraction once built.
\end{solution}

\begin{solution}{exo:g12:classes-and-objects:8}
\code{B in}, \code{A}, \code{B out}: \code{super().__init__()} runs the \omterm{def:g12:classes-and-objects:constructor}{constructor} of \code{A} in the middle of that
of \code{B}. Without the call, \code{B in} and \code{B out} only: the \omterm{def:g12:classes-and-objects:constructor}{constructor} of \code{A} never runs, and any
\omterm{def:g12:classes-and-objects:class}{attribute} it would have set is missing from the object.
\end{solution}

\begin{solution}{exo:g12:classes-and-objects:9}
\smallskip\noindent\hspace*{2em}\code{def previous_day(self):}

\noindent\hspace*{4em}\code{year, month, day = self._ymd}

\noindent\hspace*{4em}\code{if day > 1:}

\noindent\hspace*{6em}\code{return Date(year, month, day - 1)}

\noindent\hspace*{4em}\code{if month > 1:}

\noindent\hspace*{6em}\code{return Date(year, month - 1, days_in_month(year, month - 1))}

\noindent\hspace*{4em}\code{return Date(year - 1, 12, 31)}

\smallskip\noindent 1 March 2028 gives \code{2028-02-29}, a leap year; 1 January 2027 gives \code{2026-12-31}. The
tests compare the \omterm{def:g12:classes-and-objects:class}{method} with the calendar of the \omterm{def:g11:modules-and-interfaces:library}{standard library} on $1200$ days in a row.
\end{solution}

\begin{solution}{exo:g12:classes-and-objects:10}
Every fraction is made by the \omterm{def:g12:classes-and-objects:constructor}{constructor}, since \code{+} and \code{*} return \code{Fraction(...)}. So it is enough
to check the \omterm{def:g12:classes-and-objects:constructor}{constructor}, for any numerator $a$ and denominator $b \ne 0$. Let $g = \gcd(a, b)$, positive since
$b \ne 0$. The divisions $a / g$ and $b / g$ are exact, and two numbers divided by their greatest common divisor have
no common divisor other than $1$: a common divisor $d > 1$ of $a/g$ and $b/g$ would make $dg$ a common divisor of
$a$ and $b$ larger than $g$. If $b < 0$, both are divided by $-g$ instead: the signs change, the common divisors do
not, and the denominator $b / (-g)$ is positive. For \code{+} and \code{*}, the denominator passed is the product of
two positive denominators, never $0$, so the \omterm{def:g12:classes-and-objects:constructor}{constructor} accepts it and the argument applies.
\end{solution}

\begin{solution}{exo:g12:classes-and-objects:11}
\smallskip\noindent\hspace*{2em}\texttt{total = sum(map(lambda x: x * x, filter(lambda x: x \% 2 == 0, t)))}

\smallskip\noindent or, with a \omterm{def:g10:tuples-sets-comprehensions:comprehension}{comprehension}, \texttt{sum(x * x for x in t if x \% 2 == 0)}. For
\code{t = range(7)}: $0 + 4 + 16 + 36 = 56$.
\end{solution}

\begin{solution}{exo:g12:classes-and-objects:12}
\smallskip\noindent\hspace*{2em}\code{class Stack:}

\noindent\hspace*{4em}\code{def __init__(self):}

\noindent\hspace*{6em}\code{self._items = []}

\noindent\hspace*{4em}\code{def push(self, x):}

\noindent\hspace*{6em}\code{self._items.append(x)}

\noindent\hspace*{4em}\code{def pop(self):}

\noindent\hspace*{6em}\code{if not self._items:}

\noindent\hspace*{8em}\code{raise IndexError("pop from an empty stack")}

\noindent\hspace*{6em}\code{return self._items.pop()}

\noindent\hspace*{4em}\code{def peek(self):}

\noindent\hspace*{6em}\code{if not self._items:}

\noindent\hspace*{8em}\code{raise IndexError("peek at an empty stack")}

\noindent\hspace*{6em}\code{return self._items[-1]}

\noindent\hspace*{4em}\code{def is_empty(self):}

\noindent\hspace*{6em}\code{return self._items == []}

\noindent\hspace*{4em}\code{def __len__(self):}

\noindent\hspace*{6em}\code{return len(self._items)}

\smallskip\noindent Tests: after pushing $1$ then $2$, \code{pop()} returns $2$, then \code{peek()} returns $1$ and
\code{len} is $1$; a new stack is empty, and \code{pop()} on it raises \code{IndexError} (checked with
\code{assertRaises}). The next chapter builds stacks without Python's list.
\end{solution}

\begin{solution}{exo:g12:classes-and-objects:13}
Nine: \code{Fraction(0)}, then at each of the four turns \code{Fraction(1, k)} and the result of \code{+}, made by
\code{__add__}. At the end only one is reachable, the one named \code{total}, $25/12$; the eight others are no longer
named by anything, and Python frees their memory.
\end{solution}

\begin{solution}{exo:g12:classes-and-objects:14}
\code{items = []} in the body of the \omterm{def:g12:classes-and-objects:class}{class} makes one list, an \omterm{def:g12:classes-and-objects:class}{attribute} of the \omterm{def:g12:classes-and-objects:class}{class} itself. \code{self.items} looks
for an \omterm{def:g12:classes-and-objects:class}{attribute} of the object, finds none, and finds the \omterm{def:g12:classes-and-objects:class}{class}'s: every basket appends to the same list. The repair
gives each object its own list in the \omterm{def:g12:classes-and-objects:constructor}{constructor}:

\smallskip\noindent\hspace*{2em}\code{def __init__(self):}

\noindent\hspace*{4em}\code{self.items = []}
\end{solution}

\begin{solution}{exo:g12:classes-and-objects:15}
\smallskip\noindent\hspace*{2em}\code{class Polynomial:}

\noindent\hspace*{4em}\code{def __init__(self, coefficients):}

\noindent\hspace*{6em}\code{c = list(coefficients)}

\noindent\hspace*{6em}\code{while len(c) > 1 and c[-1] == 0:}

\noindent\hspace*{8em}\code{c.pop()}

\noindent\hspace*{6em}\code{self.c = c}

\noindent\hspace*{4em}\code{def __add__(self, other):}

\noindent\hspace*{6em}\code{n = max(len(self.c), len(other.c))}

\noindent\hspace*{6em}\code{a = self.c + [0] * (n - len(self.c))}

\noindent\hspace*{6em}\code{b = other.c + [0] * (n - len(other.c))}

\noindent\hspace*{6em}\code{return Polynomial([x + y for x, y in zip(a, b)])}

\noindent\hspace*{4em}\code{def __mul__(self, other):}

\noindent\hspace*{6em}\code{out = [0] * (len(self.c) + len(other.c) - 1)}

\noindent\hspace*{6em}\code{for i, x in enumerate(self.c):}

\noindent\hspace*{8em}\code{for j, y in enumerate(other.c):}

\noindent\hspace*{10em}\code{out[i + j] += x * y}

\noindent\hspace*{6em}\code{return Polynomial(out)}

\noindent\hspace*{4em}\code{def __call__(self, x):}

\noindent\hspace*{6em}\code{result = 0}

\noindent\hspace*{6em}\code{for a in reversed(self.c):}

\noindent\hspace*{8em}\code{result = result * x + a}

\noindent\hspace*{6em}\code{return result}

\smallskip\noindent The \omterm{def:g12:classes-and-objects:constructor}{constructor} drops the zero coefficients at the end, so that $(1 + x)^2 - x^2$ is stored as
\code{[1, 2]}. The term $x^i$ times $x^j$ goes to $x^{i+j}$. \code{__call__} evaluates from the highest coefficient
down, one multiplication per coefficient: $(1 + x)^2$ at $2$ gives $9$.
\end{solution}

\begin{solution}{pb:g12:classes-and-objects:1}
\textbf{1.} \code{_num} and \code{_den}. The underscore says they are internal: users go through the \omterm{def:g12:classes-and-objects:class}{methods}, which
keep the invariant.

\textbf{2.} $\gcd(10, -4) = 2$; the denominator is negative, so both are divided by $-2$: \code{_num} is $-5$ and
\code{_den} is $2$.

\textbf{3.} \code{__add__}, with \code{self} the fraction named \code{exact} and \code{other} the new
\code{Fraction(1, n)}.

\textbf{4.} The \omterm{def:g12:classes-and-objects:constructor}{constructor} stored $2/4$ as $1/2$; \code{__eq__} then compares $1$ with $1$ and $2$ with $2$.

\textbf{5.} The \omterm{def:g12:classes-and-objects:class}{methods} give read access only, so that no user can break the invariant by writing; and the \omterm{def:g12:classes-and-objects:class}{class}
stays free to store its data otherwise without changing its users.

\textbf{6.} The denominator is positive, and the numerator and denominator have no common divisor other than $1$.

\textbf{7.} Exercise \ref{exo:g12:classes-and-objects:10}.

\textbf{8.} Let $a/b$ and $c/d$ be reduced, with $b, d > 0$. If they are equal numbers, $ad = bc$. Then $b$ divides
$ad$ and has no common divisor with $a$, so $b$ divides $d$; in the same way $d$ divides $b$. Both being positive,
$b = d$, and then $a = c$. Equal numbers have equal \omterm{def:g12:classes-and-objects:class}{attributes}, so comparing \omterm{def:g12:classes-and-objects:class}{attributes} is right; without the
invariant, $1/2$ and $2/4$ would be found different.

\textbf{9.} That it does not protect them: the underscore is a convention, and any program can write
\code{half._num = 2}. The invariant breaks, and \code{__eq__}, which relies on it, answers wrongly that $2/2$ is not
$1$.

\textbf{10.} $\frac{12}{12} + \frac{6}{12} + \frac{4}{12} + \frac{3}{12} = \frac{25}{12}$, as printed.

\textbf{11.} $H_9 = 7129/2520$ and $H_{10} = 7381/2520$; $\frac{1}{10} = \frac{252}{2520}$, and $7129 + 252 = 7381$.

\textbf{12.} $434$ digits. A $64$-bit integer holds at most $2^{64} - 1$, which has $20$ digits; Python's integers grow
as needed (chapter \ref{ch:g10:machine-integers}).

\textbf{13.} A \omterm{def:g10:floating-point:float}{floating-point number} is a whole number of at most $53$ bits times a power of $2$
(chapter \ref{ch:g10:floating-point}): its denominator is a power of $2$. $11/6$ has the factor $3$ in its reduced
denominator, so its binary writing never ends, and every \omterm{def:g10:floating-point:float}{floating-point number} differs from it.

\textbf{14.} Each $k$ from $1$ to $n$ divides $\operatorname{lcm}(1, \dots, n)$, so $\frac1k$ is the fraction with
that denominator and numerator $\operatorname{lcm}(1, \dots, n) / k$, a whole number.

\textbf{15.} For $n = 6$, $7$ and $8$, by a factor $3$; for $n = 20$, by a factor $15$. For the other printed values
the reduced denominator is the common multiple itself.

\textbf{16.} $\operatorname{lcm}(1, \dots, 6) = 60$, and $H_6 = \frac{60 + 30 + 20 + 15 + 12 + 10}{60} =
\frac{147}{60}$. $147 = 3 \times 49$ and $60 = 3 \times 20$: $H_6 = 49/20$.

\textbf{17.} $0$ at $n = 10$ and $n = 100$: the floating-point sum is the closest \omterm{def:g10:floating-point:float}{floating-point number} to the exact
value. At $n = 1000$, $-1.8 \times 10^{-15}$.

\textbf{18.} Each addition rounds its result to the nearest \omterm{def:g10:floating-point:float}{floating-point number}, and $1/k$ is itself rounded; the
small errors of a thousand roundings add up instead of cancelling out.

\textbf{19.} The value is about $7.49$, so the error is about $2.4 \times 10^{-16}$ of it: two units of the last bit
for numbers between $4$ and $8$, whose spacing is $2^{-50} \approx 8.9 \times 10^{-16}$. Fifteen significant digits
are right.

\textbf{20.} \textbf{$H_{10} = 7381/2520$ exactly; summed in floating point up to $n = 1000$, the harmonic sum is off
by $1.8 \times 10^{-15}$, about $2 \times 10^{-16}$ of its value $7.485\,470\,860\,550\,345$.}
\end{solution}
